\section{A comprehensive description of the social planner problem}\label{sec:app:sp}

This appendix carries the full characterisation of Problem~\ref{prob:sp}. Section~\ref{sec:sp:opt} states the conditions compactly; what follows derives each of them, decomposes it into the separate margins it balances, and checks the units. The units are worth taking seriously in this model, because two distinct quantities are being traded off throughout --- final goods and tonnes of material --- and most of the conditions below mix them.

\subsection{Units}\label{sec:app:sp:units}

Everything in the model is measured in one of two units, the final good and the tonne, or in a ratio of the two. Collecting them once:
\begin{itemize}
\item \emph{Final goods.} Output $Y$, consumption $C$, investment $I$, the capital stock $K$, gross formation $G$, the cost flows $C^N$, $C^D$, $C^W$, and the waste base $\mathcal D$.
\item \emph{Tonnes.} Raw material $R$ and its components $N$ and $R^R$, waste generation $W$, the waste stock $\mathcal W$ and the handled flow $H=\mu\mathcal W$, treated waste $T=\varpi H$, the reserve $S$, cumulative discoveries $X$ and their ceiling $\bar X_{\max}$, the discovery flow $D$, embodied material $M^K$, the pollution stock $P$, the emission flow $\Xi$, and the nature-attributable base $\mathcal N$. Flows are tonnes per period; with the period as the unit of time, flows and stocks are directly comparable.
\item \emph{Tonnes per final good.} The material intensities $\Omega$, $\Omega^j=\Omega\omega^j$, $\Phi^I=\Omega\phi^I$, $\Phi^K$ and $\Phi^Y$; also $a'\left(x\right)$, the marginal product of recycling capital.
\item \emph{Final goods per tonne.} The marginal product of materials $F_R$, the marginal costs $C^N_N$, $C^D_D$, $c^W$ and $c^{T\prime}$, the recycling capital ratio $x=K^R/T$, the handling dividend $h$, and every shadow price carried below: $p^S$, $p^X$, $p^W$, $p^{\mathcal W}$, $p^M$, $p^P$, $\zeta$ and $\Psi$.
\item \emph{Dimensionless.} The relative intensities $\omega^j$, the durables coefficient $\phi^I$, the recycling functions $a$, $\alpha$ and $\varepsilon$ and the yield ceiling $\bar a$, the leakage rate $d^W$, the treatment share $\varpi$, the handling share $\mu$ (the fraction of the stock handled per period), the storage share $\sigma$, the marginal product of capital $F_K$, the shadow price of capital $q$, and the nature-attributable intensities $\Omega^{N,S}$ and $\Omega^{D,S}$.
\end{itemize}
Two of these deserve emphasis. The intensities $\Omega^j$ convert a \emph{final-good} flow into tonnes and so carry units, whereas $\Omega^{N,S}$ and $\Omega^{D,S}$ convert a tonnage into a tonnage and are pure numbers; the two families are not comparable, and the difference is exactly the distinction of Section~\ref{sec:sp:setup:materialaccounting} between mass drawn through $R$ and mass displaced from nature. And $\lambda$, the multiplier on the goods constraint, is measured in utils per final good, which is why dividing by it puts every shadow price into final-good units as in~\eqref{eq:sp:costates-def}.

A useful consequence: any product of a shadow price with an intensity, such as $\zeta\Omega^C$ or $\Phi^Ip^W$, is dimensionless and can be read directly as a proportional wedge.

One notational convention is carried over from Section~\ref{sec:sp:setup:primitives}: the primitives $F$, $a$, $C^N$, $C^D$, $c^c$ and $c^T$ all shift exogenously with calendar time, and the argument $t$ is suppressed throughout this appendix, as is the period subscript wherever a condition is a within-period statement. Nothing below depends on the suppression. Since $t$ is not a choice variable it contributes no first-order condition of its own, and every condition holds period by period with the primitives evaluated at that date. Technological change matters for how the system evolves --- the problem is non-autonomous --- not for the form of the conditions that characterise it.

\subsection{First-order conditions for the controls}\label{sec:app:sp:foc}

Each condition below is obtained by differentiating the Lagrangian~\eqref{eq:sp:lagrangian} with respect to a period-$t$ control, and each is presented the same way: what the control moves, the condition itself with its terms marked, and then the terms one at a time. Every control acts on flows and constraints of its own period alone --- the sole intertemporal consequence of a control, the stock its period-$t$ choice feeds, is priced by the costate dated $t$ under the convention of~\eqref{eq:sp:costates-def} --- so each derivative involves period-$t$ terms only.

\smalltitle{Material intensity.}
Differentiation of~\eqref{eq:sp:lagrangian} with respect to $\Omega$ gives
\begin{align*}
    \frac{\partial \mathcal{L}}{\partial \Omega} = \lambda \phi^IG\left(p^W-p^M\right) - \lambda \zeta \left(\mathcal{D}+\phi^IG\right),
\end{align*}
so that the optimum is characterised by
\begin{align}
    \underbrace{\phi^IG}_{\text{(a)}}\underbrace{\left(p^W-p^M\right)}_{\text{(b)}}
    &= \underbrace{\zeta}_{\text{(c)}}\underbrace{\left(\mathcal D+\phi^IG\right)}_{\text{(d)}}
    \;=\;\zeta\,\frac{R}{\Omega}.
\label{eq:sp:foc:Omega}
\end{align}
Each piece has a direct reading.

Term (a) is the tonnage that a marginal rise in $\Omega$ diverts into new capital: with $\Phi^I=\Omega\phi^I$ and gross formation $G$, a perturbation $d\Omega$ raises embodied investment by $\phi^IG\,d\Omega$ tonnes. This is the only \emph{physical} margin that $\Omega$ moves. It is worth being clear about what it does not move: given $R$, the ledger~\eqref{eq:sp:ledger} states that every tonne entering production becomes waste except what is stored in the capital stock, so $\Omega$ shifts the split between the two and never the total.

Term (b) prices that shift, in final goods per tonne. The diverted tonne leaves the current waste flow, saving $p^W$, and enters the embodied stock $M^K$, incurring $p^M$. It is a postponement, not a disposal, which is why the two shadow prices enter as a difference.

Term (d) is the embodied-material requirement of the entire use structure, $\mathcal D+\phi^IG=R/\Omega$ by the intensity condition~\eqref{eq:sp:intensity}: the perturbation $d\Omega$ raises the material that the final-good system must carry by $\left(\mathcal D+\phi^IG\right)d\Omega$ tonnes. Since $R$ is held fixed along the perturbation, the increase has to be paid for, and term (c) is the price. Thus $\zeta$ is the shadow value, in final goods per tonne, of one tonne of embodied material in the final-good system --- what it would be worth to have that tonne appear without drawing it through $R$.

Both sides are in final goods per unit of $\Omega$: (a) and (d) are in final goods, (b) and (c) in final goods per tonne.

\smalltitle{Consumption.} A marginal unit of consumption is one final good. It raises felicity by $u'\left(C\right)$ utils, absorbs one unit of the goods constraint, and adds $\omega^C$ to the waste base $\mathcal D$, hence $\Omega^C$ tonnes to the embodied material the system must carry. Differentiating and dividing by $\lambda$ to put the condition in final goods,
\begin{align}
\underbrace{\frac{u'\left(C\right)}{\lambda}}_{\text{(a)}} &= \underbrace{1}_{\text{(b)}}+\underbrace{\zeta\Omega^C}_{\text{(c)}} .
\label{eq:sp:foc:C}
\end{align}
Term (a) is the marginal utility of consumption converted into final goods by the goods multiplier, and is dimensionless. Term (b) is the direct resource cost, one good for one good. Term (c) is the material wedge: $\Omega^C$ tonnes per good times $\zeta$ goods per tonne, again dimensionless, and readable directly as a proportional tax.

The condition says that $\lambda$ is \emph{not} the marginal utility of consumption in this model, as it would be without the materials accounting. A unit consumed also carries embodied material away from the storage margin, and $\zeta\Omega^C$ is what that costs. Because every shadow price in~\eqref{eq:sp:costates-def} is normalised by $\lambda$ rather than by $u'\left(C\right)$, the distinction matters for reading all of them. The wedge is a tax when $\zeta>0$ and a subsidy when $\zeta<0$; by~\eqref{eq:sp:zeta-explicit} its sign is that of $p^W-p^M$, never that of $p^W$ alone.

\smalltitle{Investment.} A marginal unit of investment is one final good. It absorbs one unit of the goods constraint, adds one unit to next period's capital stock, and --- since $G=I+\delta K$ --- raises embodied investment $\Omega\phi^IG$ by $\Phi^I$ tonnes, which appears in three places at once: the ledger, the transition for $M^K$, and the right-hand side of the intensity condition. It does \emph{not} enter $\mathcal D$, because the material embodied in capital is carried by the separate term $\phi^IG$ in~\eqref{eq:sp:setup:ybalance_OmegaPhi}. Hence $\partial\mathcal L/\partial I=\lambda\left[-1+q+\Omega\phi^I\left(p^W-p^M-\zeta\right)\right]$ and
\begin{align}
\underbrace{q}_{\text{(a)}} &= \underbrace{1}_{\text{(b)}}-\underbrace{\Phi^I}_{\text{(c)}}\Big[\underbrace{\left(p^W-p^M\right)}_{\text{(d)}}-\underbrace{\zeta}_{\text{(e)}}\Big]
\;=\;1-\Phi^I\left(1-\sigma\right)\left(p^W-p^M\right).
\label{eq:sp:foc:I}
\end{align}
Term (a) is the shadow price of installed capital, dimensionless because $K$ is measured in final goods. Term (b) is the resource cost of the unit. Term (c) is the tonnage stored per good invested. Terms (d) and (e) are both in final goods per tonne, so the product with (c) is dimensionless and comparable with (a) and (b).

Terms (d) and (e) pull in opposite directions and it is worth separating them. Term (d) is the gross storage benefit already met in~\eqref{eq:sp:foc:Omega}: the tonne avoids today's waste flow and joins the embodied stock instead. Term (e) is the claw-back. Embodying material in capital also \emph{consumes} embodied material, which must be drawn through $R$ and is therefore charged at $\zeta$; by~\eqref{eq:sp:zeta-explicit} this recovers exactly the fraction $\sigma$ of the gross benefit, leaving a net storage premium of $\left(1-\sigma\right)\left(p^W-p^M\right)$ per tonne.

The reading is that $q$ departs from unity by precisely the value of the material a unit of gross investment stores. When $p^W>p^M$ the departure is downward: the planner is willing to invest even where the capital services alone would not justify it, because investing also holds mass out of the current waste flow.

\smalltitle{The value of raw material.} The four remaining conditions all price a tonne delivered to production, and the same combination recurs in each. Let
\begin{align}
\Psi &\equiv \underbrace{F_R\left(1-\zeta\Omega^Y\right)}_{\text{output, net of its own material}}+\underbrace{\zeta}_{\text{relaxation of the intensity condition}}
\label{eq:sp:Psi}
\end{align}
be the \emph{gross material value}, in final goods per tonne. The first term is the output gained from one more tonne of $R$, scaled down because that extra output itself requires $\Omega^Y$ tonnes per good of embodied material, charged at $\zeta$; the factor $1-\zeta\Omega^Y$ is dimensionless. The second term is the direct effect of the tonne on the intensity condition~\eqref{eq:sp:intensity}, where $R$ appears on the left-hand side.

What $\Psi$ deliberately omits is the ledger. The tonne will eventually leave the economy as waste, and that consequence is priced by $p^W$, which appears separately in each condition below. Keeping the two apart is what makes the conditions readable: $\Psi$ is what a tonne is worth on its way in, $p^W$ what it costs on its way out.

\smalltitle{Extraction.} A marginal tonne extracted adds one tonne to $R$, costs $C^N_N$ final goods of extraction effort, draws the reserve down by one tonne, and displaces $\Omega^{N,S}$ tonnes of overburden that never pass through $R$. Collecting the five terms of $\partial\mathcal L/\partial N$,
\begin{align}
\underbrace{\Psi}_{\text{(a)}} &= \underbrace{C^N_N\left(1+\zeta\Omega^N\right)}_{\text{(b)}}
+\underbrace{p^S}_{\text{(c)}}
+\underbrace{p^W\left(1+\Omega^{N,S}\right)}_{\text{(d)}} .
\label{eq:sp:foc:N}
\end{align}
Every term is in final goods per tonne.

Term (a) is the value of the tonne delivered, from~\eqref{eq:sp:Psi}. Term (b) is the extraction effort, grossed up because the effort is itself final-good expenditure and therefore carries $\Omega^N$ tonnes per good of embodied material; the factor $1+\zeta\Omega^N$ is the same kind of dimensionless wedge as in~\eqref{eq:sp:foc:C}. Term (c) is the in-situ rent: extracting now forecloses extracting later, priced by the reserve costate. Term (d) charges the tonne twice over. The $1$ is the ledger: every tonne entering production must eventually leave as waste. The $\Omega^{N,S}$ is the displaced mass, which becomes waste without ever having been material input; it is dimensionless, so the product with $p^W$ is again final goods per tonne.

\smalltitle{Exploration.} A marginal tonne of discovery costs $C^D_D$ final goods, raises both $S$ and $X$ by one tonne, and displaces $\Omega^{D,S}$ tonnes. It adds nothing to $R$. Hence
\begin{align}
\underbrace{p^S}_{\text{(a)}}+\underbrace{p^X}_{\text{(b)}} &= \underbrace{C^D_D\left(1+\zeta\Omega^D\right)}_{\text{(c)}}+\underbrace{p^W\Omega^{D,S}}_{\text{(d)}} .
\label{eq:sp:foc:D}
\end{align}
Term (a) is the rent created by the discovery and term (b) the shadow cost of having used up an easy deposit, with $p^X<0$; their sum is the net value of a discovery. Term (c) is the grossed-up effort, exactly as in~\eqref{eq:sp:foc:N}. Term (d) is the waste from displaced mass alone.

The comparison with~\eqref{eq:sp:foc:N} is the informative one. Both activities move tonnes and both are grossed up for the embodied material in their effort, but the unit charge $p^W$ present in extraction is absent here, because exploration adds nothing to $R$ and therefore generates no waste beyond what it displaces. Discovery is not a substitute for extraction in the materials ledger; it only relaxes the constraint that makes extraction expensive.

Exhaustibility acts on this margin, and only on this margin. By~\eqref{eq:sp:setup:exhaustible-discovery} the marginal cost of the first tonne, $C^D_D\left(0,X,t\right)$, diverges as $X\to\bar X_{\max}$, while the left-hand side of~\eqref{eq:sp:foc:D} stays bounded: $p^X\le0$, and $p^S$ is finite along any admissible path. So~\eqref{eq:sp:foc:D} can hold with $D>0$ only while $X$ is strictly below its ceiling; exploration must eventually cease or dwindle to nothing, and $X$ approaches $\bar X_{\max}$ at most asymptotically. No constraint $X\le\bar X_{\max}$ ever has to be adjoined to the problem --- the cost function polices the ceiling by itself --- and cumulative virgin extraction inherits the bound, $\sum_{t=0}^\infty N_t\le S_0+\bar X_{\max}-X_0$, as recorded at~\eqref{eq:sp:sum:cumulative-N}.

\smalltitle{Waste.} This is the condition that identifies $p^W$, and with the waste stock it is the simplest condition in the system. In the period of generation a marginal tonne of waste does exactly one thing: it enters next period's stockpile. Collection and treatment, recycling, and emissions are all levied on the handled flow $H=\mu\mathcal W$, which is a function of the predetermined state rather than of the control, so $W_t$ appears in the Lagrangian~\eqref{eq:sp:lagrangian} in two places only --- the ledger term $+\lambda_t p^W_tW_t$ and the stock's transition $+\lambda_t p^{\mathcal W}_tW_t$ --- and
\begin{align}
\underbrace{p^W_t}_{\text{(a)}} &= \underbrace{-p^{\mathcal W}_t}_{\text{(b)}} .
\label{eq:sp:foc:W}
\end{align}
Both sides are in final goods per tonne. Term (a) is the social cost the ledger attributes to generating a tonne; term (b) says that this cost is nothing but the negative of the value of the tonne where it lands. The identification is contemporaneous and carries no discount factor, because $p^{\mathcal W}_t$ is defined under~\eqref{eq:sp:costates-def} as the period-$t$ value of a tonne that enters the stock at $t+1$ --- exactly the tonne being generated. Generating waste is neither taxed for its handling nor credited for its feedstock at this margin --- both lie in the future, in the periods the stockpile releases the tonne --- and they enter through the costate recursion~\eqref{eq:sp:costate:Wstock} below, which prices the stockpile as an asset.

The sign of $p^W$ therefore still follows from a comparison, not an assumption, but the comparison is forward-looking: waste is a social cost when the discounted handling dividends of a stockpiled tonne are negative --- handling resources and emission damages outweighing the feedstock recovered --- and a social \emph{asset} in the reverse case, the configuration of squeeze paths on which the stockpile is the economy's remaining source of material supply.

\smalltitle{Treatment share.} The treatment share is dimensionless, so a perturbation $d\varpi$ routes a further $H\,d\varpi$ tonnes to treatment, of which $\alpha H\,d\varpi$ becomes secondary material, at marginal cost $\left[c^c+c^{T\prime}\left(\varpi\right)\right]H\,d\varpi$ final goods. Dividing $\partial\mathcal L/\partial\varpi$ through by $\lambda H$ to state the condition per handled tonne, the marginal gain is
\begin{align}
\mathcal G^\varpi \;\equiv\;
\underbrace{\alpha}_{\text{(a)}}\Big[\underbrace{\Psi}_{\text{(b)}}-\underbrace{p^W}_{\text{(c)}}\Big]
+\underbrace{p^P\left[\left(1-d^W\right)+d^W\alpha\right]}_{\text{(d)}} ,
\label{eq:sp:foc:varpi-gain}
\end{align}
to be set against the marginal cost $\underbrace{\left[c^c+c^{T\prime}\left(\varpi\right)\right]\left(1+\zeta\Omega^W\right)}_{\text{(e)}}$. All terms are in final goods per handled tonne.

Term (a) is the marginal recycling yield $\alpha=a-a'x$, dimensionless: the tonnes of secondary material obtained from one more tonne of treated waste, holding recycling capital fixed. It is the marginal and not the average yield $a$ that appears, because treating more waste with a given $K^R$ dilutes the capital across a larger throughput.

Term (b) values that material as in~\eqref{eq:sp:Psi}. Term (c) is the part worth dwelling on, and it is a direct consequence of the ledger. Recycling a tonne does not remove it from the economy; it returns it to $R$, from which it must leave as waste all over again. Recycling therefore \emph{raises} total waste generation even as it lowers extraction --- this is the circularity multiplier of~\eqref{eq:sp:setup:ledger-collected} seen from the margin --- so the material recovered is worth $\Psi-p^W$, not $\Psi$.

Term (d) collects two distinct pollution gains. The treated tonne no longer joins the untreated stream, worth $1-d^W$ per tonne; and the fraction recovered as material does not leak either, worth a further $d^W\alpha$. Term (e) is the marginal handling cost --- the collection charge $c^c$ on the marginal collected tonne plus the treatment cost, rising in $\varpi$ by assumption~\eqref{eq:sp:setup:waste-cost} --- grossed up for its own embodied material; the convexity of $c^T$ is what makes the interior branch below a single well-defined root rather than a set.

\smalltitle{Treatment share: the bounds.} Because $\varpi$ is bounded on both sides, this margin is
not an equality in general. Attach multipliers $\lambda\nu^0\ge0$ to $\varpi\ge0$ and
$\lambda\nu^1\ge0$ to $\varpi\le1$. Then $\partial\mathcal L/\partial\varpi=0$ reads
\begin{align}
\mathcal G^\varpi-\left[c^c+c^{T\prime}\left(\varpi\right)\right]\left(1+\zeta\Omega^W\right)
&= \nu^1-\nu^0,
&
\nu^0\varpi &= 0,
&
\nu^1\left(1-\varpi\right) &= 0,
\label{eq:sp:foc:varpi-kkt}
\end{align}
which is~\eqref{eq:sp:sum:varpi} of the main text once the two complementary-slackness conditions are
resolved into branches: $\nu^0>0$ gives $\varpi=0$ and
$\mathcal G^\varpi\le c^c\left(1+\zeta\Omega^W\right)$; $\nu^1>0$ gives
$\varpi=1$ and $\mathcal G^\varpi\ge \left[c^c+c^{T\prime}\left(1\right)\right]\left(1+\zeta\Omega^W\right)$;
$\nu^0=\nu^1=0$ gives the interior equality, which by convexity of $c^T$ has a root in
$\left(0,1\right)$ exactly when $\mathcal G^\varpi$ lies strictly between the two bounding marginal
costs.

Two features of the branch structure are worth marking. First,
$c^{T\prime}\left(0\right)=0$ by~\eqref{eq:sp:setup:waste-cost}, so the marginal cost of the first
treated tonne is the collection charge $c^c$ alone, and the lower branch is a genuine choke: the
economy dumps its handled waste untreated whenever the gain cannot cover collection. The corner is
economically meaningful --- it is the configuration of an economy whose output has collapsed and
which can no longer pay to process the stockpile draining past it --- and the collection charge is
what keeps it feasible. One caveat survives: at $\varpi=0$ the ratio $x=K^R/\left(\varpi H\right)$ is
undefined and $\alpha$ is not pinned down there, so entry into treatment and entry into recycling
capital are one decision at the corner --- the corner test should be run jointly with the capital
margin, over the best available $\left(\varpi,x\right)$ deviation, and the numerical implementation
should treat it as such.
The upper branch carries no such degeneracy: $c^c+c^{T\prime}\left(1\right)$ is a finite number
under~\eqref{eq:sp:setup:waste-cost}, so $\varpi=1$ is a genuine possibility, ruled neither in nor out
by the assumptions made so far, and $\nu^1>0$ is a case the numerical implementation has to test for.

\smalltitle{Recycling capital.} A marginal unit of recycling capital is one final good moved out of final-goods production. It costs $F_K$ of output and yields $a'$ tonnes of secondary material. Writing the marginal gain per tonne recovered as
\begin{align}
\mathcal G^K \;\equiv\; \underbrace{\Psi-p^W}_{\text{(b)}}+\underbrace{d^Wp^P}_{\text{(c)}},
&&\text{the interior condition is}
&&
\underbrace{a'}_{\text{(a)}}\,\mathcal G^K &= \underbrace{F_K\left(1-\zeta\Omega^Y\right)}_{\text{(d)}} .
\label{eq:sp:foc:KR}
\end{align}
Every term is dimensionless: (a) is tonnes per final good, $\mathcal G^K$ is final goods per tonne, and (d) is a marginal product of capital adjusted by a dimensionless factor.

Term (a) is the marginal product of recycling capital, $\mathcal R_{K^R}=a'$. Term (b) values the tonne recovered exactly as in~\eqref{eq:sp:foc:varpi-gain}, net of the waste it will regenerate. Term (c) is the pollution gain, and the contrast with term (d) of~\eqref{eq:sp:foc:varpi-gain} is instructive: recycling capital improves the yield on waste that is \emph{already} being treated, so it cannot claim the $\left(1-d^W\right)$ diversion from the untreated stream and collects only the $d^W$ leakage avoided on the extra material recovered. Term (d) is the opportunity cost, itself net of embodied material, because output forgone also releases the material that output would have required.

\smalltitle{Recycling capital: the bound.} Attaching a multiplier $\lambda\nu^K\ge0$ to $K^R\ge0$,
$\partial\mathcal L/\partial K^R=0$ reads
\begin{align}
a'\!\left(x\right)\mathcal G^K-F_K\left(1-\zeta\Omega^Y\right) &= -\nu^K,
&
\nu^KK^R &= 0 .
\label{eq:sp:foc:KR-kkt}
\end{align}
Only the lower bound is attached. The upper bound $K^R\le K$ implied by $K^Y\ge0$ is
slack at any candidate optimum under the maintained Inada condition on $K^Y$, since $F_K\to\infty$ as
$K^Y\to0$ makes the right-hand side of~\eqref{eq:sp:foc:KR} unboundedly large against a finite left-hand
side.

Resolving~\eqref{eq:sp:foc:KR-kkt} into branches gives~\eqref{eq:sp:sum:KR} of the main text. What
makes the resolution clean --- and this is the difference from the treatment margin --- is that
$a'$ is strictly decreasing with a finite value at the origin. The marginal gain
$a'\!\left(x\right)\mathcal G^K$ is therefore maximised at $x=0$ and falls monotonically thereafter,
so the single number $a'\!\left(0^+\right)\mathcal G^K$ decides between the two branches: if it does
not clear $F_K\left(1-\zeta\Omega^Y\right)$ no positive $x$ will either and the corner is optimal,
and if it does clear it there is exactly one interior root. Note that the comparison is between
endogenous objects evaluated at the candidate allocation, with $F_K$ read at $K^Y=K$ in the corner
branch, so the branching inequality is not a restriction on primitives.

Read together, \eqref{eq:sp:foc:varpi-kkt} and~\eqref{eq:sp:foc:KR-kkt} show how the two arguments of the recycling technology are priced separately: the treatment margin by the marginal yield $\alpha=a-a'x$, the capital margin by $a'$. The constant-returns form~\eqref{eq:sp:setup:recycling-crs} guarantees that the two exhaust the product, $\mathcal R=\alpha T+a'K^R$.

\smalltitle{The avoided cost of virgin material.} Both recycling conditions are stated above in terms
of $\Psi-p^W$, the value of a recovered tonne net of the waste it will regenerate. That form is the
right one for deriving them, but it is not the clearest one for reading them, because $\Psi$ and
$p^W$ are both endogenous prices whose signs and magnitudes are themselves outcomes. Whenever
extraction is interior, $N>0$, the pair can be eliminated. Solving the extraction
condition~\eqref{eq:sp:foc:N} for $\Psi$ and substituting, the unit charge $p^W$ cancels save for the
overburden term:
\begin{align}
\Psi-p^W &= \underbrace{C^N_N\left(1+\zeta\Omega^N\right)}_{\text{extraction effort}}
+\underbrace{p^S}_{\text{in-situ rent}}
+\underbrace{\Omega^{N,S}p^W}_{\text{overburden}} ,
&&\text{hence}
&
V &\equiv \Psi-p^W+d^Wp^P ,
\label{eq:sp:virgin-value}
\end{align}
where we name the combination
\begin{align}
V &= C^N_N\left(1+\zeta\Omega^N\right)+p^S+\Omega^{N,S}p^W+d^Wp^P
\label{eq:sp:V}
\end{align}
the \emph{virgin displacement value}, in final goods per tonne. In these terms the two marginal gains
of~\eqref{eq:sp:foc:varpi-gain} and~\eqref{eq:sp:foc:KR} become
\begin{align}
\mathcal G^K &= V,
&
\mathcal G^\varpi &= \alpha V+p^P\left(1-d^W\right),
\label{eq:sp:foc:recycling-virgin}
\end{align}
so that the Kuhn--Tucker systems~\eqref{eq:sp:foc:varpi-kkt} and~\eqref{eq:sp:foc:KR-kkt} take the
form reported at~\eqref{eq:sp:sum:recycling-virgin} of the main text, with interior conditions
$\alpha V+p^P\left(1-d^W\right)=\left[c^c+c^{T\prime}\left(\varpi\right)\right]\left(1+\zeta\Omega^W\right)$ and
$a'V=F_K\left(1-\zeta\Omega^Y\right)$.
The first of~\eqref{eq:sp:foc:recycling-virgin} is immediate. The second is the substitution
of~\eqref{eq:sp:virgin-value} into~\eqref{eq:sp:foc:varpi-gain}, and it simplifies more than one might
expect: the leakage term $\alpha d^Wp^P$ that term (d) of~\eqref{eq:sp:foc:varpi-gain} carried separately
is precisely the $d^Wp^P$ inside $V$ applied to the $\alpha$ tonnes recovered, so the two coincide
and $V$ absorbs it. What remains outside is $p^P\left(1-d^W\right)$: the diversion of the
\emph{whole} tonne from the untreated stream, which is a pollution gain and not a materials gain, and
is therefore the one part of the treatment motive that $V$ cannot price.

The reading is the one the ledger has been pointing at throughout: \emph{recycling is paid the cost
of the virgin tonne it displaces, and nothing else.} Every component of $V$ is a cost avoided at the
mine --- the extraction effort together with the embodied material that effort carries, the in-situ
rent on the reserve left in the ground, the disposal of the overburden that would have come up
alongside the ore, and the leakage $d^W$ that tonne would have suffered. What is conspicuously
absent is $p^W$ itself. It cancels because the recovered tonne and the virgin tonne it replaces face
exactly the same charge on the way out: recycling changes the \emph{route} a tonne takes into
production, never the fact that it must eventually leave. This is the same principle noted after
\eqref{eq:sp:foc:D}, seen from the other side.


\smalltitle{Signs.} It is convenient to rearrange~\eqref{eq:sp:foc:Omega}. Define the \emph{storage share} $\sigma$, the dimensionless fraction of the material entering production that is embodied in new capital rather than diffused as waste:
\begin{align}
\sigma &\equiv \frac{\Phi^IG}{R} \;=\; \frac{\phi^IG}{\mathcal D+\phi^IG}\;\in\left[0,1\right],
&&\text{so that}
&
\zeta &= \sigma\left(p^W-p^M\right).
\label{eq:sp:zeta-explicit}
\end{align}
Two things follow at once. First, $\zeta>0$ if and only if $p^W>p^M$: embodied material is valuable exactly when storing a tonne is worth more than releasing it. Second, $\zeta\to0$ as $\phi^I\to0$. With nowhere to store material, embodied material has no shadow price of its own --- conditional on $R$, the cost of the mass that eventually becomes waste is carried in full by $p^W$ through the ledger, and all that $\Omega$ prices is the timing of its release --- and the composition of final-good uses loses any allocative role.

What does \emph{not} follow is an unconditional sign for either price. The sign of $p^W$ is an outcome of the stockpile valuation~\eqref{eq:sp:costate:Wstock} through the identification~\eqref{eq:sp:foc:W}, as set out above. Nor does $p^M$ simply inherit whichever sign $p^W$ currently has: the costate recursion~\eqref{eq:sp:costate:M} below makes it a \emph{forward} average of $p^W$, not a function of its current value, so the three signs need not move together along a transition.

\subsection{Costate equations}\label{sec:app:sp:costates}

Each state $Z$ carries the costate $m$ named in~\eqref{eq:sp:costates-def}, attached in period $t$ to the transition that determines $Z_{t+1}$, together with the transversality condition $\lim_{t\to\infty}\beta^t\lambda_t m_t Z_{t+1}=0$. (We write the generic costate as $m$ rather than the customary $\mu$, which now denotes the handling share of the waste stock.) The costate condition is obtained by differentiating the Lagrangian~\eqref{eq:sp:lagrangian} with respect to $Z_{t+1}$, which appears exactly twice: in the period-$t$ terms through the transition being adjoined, contributing $-\beta^t\lambda_tm_t$, and in the period-$t{+}1$ terms through every dividend the stock pays and its own onward transition. Setting the derivative to zero and writing $1+r_{t+1}\equiv\lambda_t/\left(\beta\lambda_{t+1}\right)$ for the social discount rate implied by the goods multiplier, every condition below takes the form
\begin{align*}
\left(1+r_{t+1}\right)m_t &= \text{dividend}_{t+1}+\text{surviving fraction}\times m_{t+1},
\end{align*}
whose forward solution, selected by transversality, prices the stock as the discounted sum of the dividends it pays while it survives. Throughout, the dividend and the surviving fraction are evaluated at $t+1$ and the derivations report $\partial\mathcal L^{t+1}/\partial Z_{t+1}$, the derivative of the period-$t{+}1$ terms.

\smalltitle{Capital.} A marginal unit of next period's capital stock is one final good. In period $t+1$ it raises output by $F_K$, requires replacement investment of $\delta$ final goods --- which, through $G=I+\delta K$, embodies a further $\Phi^I\delta$ tonnes --- and survives into the onward transition with the price $q_{t+1}$ attached to its full unit. Collecting,
\begin{align*}
\frac{\partial\mathcal L^{t+1}}{\partial K_{t+1}} &= \lambda_{t+1}\Big[\underbrace{F_K\left(1-\zeta\Omega^Y\right)}_{\text{output, net of its material}}-\underbrace{\delta}_{\text{replacement}}+\underbrace{\Phi^I\delta\left(p^W-p^M-\zeta\right)}_{\text{material stored by replacement}}+\underbrace{q_{t+1}}_{\text{continuation}}\Big].
\end{align*}
The third term equals $\delta\left(1-q_{t+1}\right)$ by the investment condition~\eqref{eq:sp:foc:I}, so the entire materials block collapses into $q$ and
\begin{align}
\left(1+r_{t+1}\right)q_t &= \underbrace{F_K\left(1-\zeta\Omega^Y\right)\Big|_{t+1}}_{\text{(a)}}+\underbrace{\left(1-\delta\right)q_{t+1}}_{\text{(b)}} .
\label{eq:sp:costate:K}
\end{align}
This is the usual one-period arbitrage condition with two modifications, both confined to places where they can be read off. In term (a) the marginal product is net of the embodied material the extra output requires. Term (b) is the depreciated continuation value, unchanged from the standard model.

That the storage motive appears nowhere else is not a coincidence. Depreciation destroys capital and releases its embodied material at the same rate $\delta$, so the material liability of a unit of capital scales exactly with the premium at which it was installed, and the two cancel in the return.

\smalltitle{Reserves.} A marginal tonne of known reserves enters period $t+1$ only through the extraction cost, lowering it by $\left|C^N_S\right|$ final goods per tonne extracted, since $C^N_S<0$, and it survives in full. Hence $\partial\mathcal L^{t+1}/\partial S_{t+1}=\lambda_{t+1}\left[-C^N_S\left(1+\zeta\Omega^N\right)+p^S_{t+1}\right]$ and
\begin{align}
\underbrace{\left(1+r_{t+1}\right)p^S_t}_{\text{(a)}} &= \underbrace{\left|C^N_S\right|\left(1+\zeta\Omega^N\right)\Big|_{t+1}}_{\text{(b)}} + \underbrace{p^S_{t+1}}_{\text{(c)}} .
\label{eq:sp:costate:S}
\end{align}
Term (a) is the required gross return on a tonne held in the ground. Term (b) is the dividend the stock pays while held --- a larger reserve makes the extraction still to come cheaper --- and term (c) is the undiminished continuation value. This is Hotelling's rule with a stock effect: the rent need not grow at the rate of interest, because holding reserves yields a flow benefit as well as a capital gain, and the dividend is grossed up by $1+\zeta\Omega^N$ for the embodied material carried by extraction effort.

\smalltitle{Discoveries.} A marginal tonne of cumulative discovery raises the cost of period-$t{+}1$ exploration by $C^D_X>0$, and does nothing else. So
\begin{align}
\left(1+r_{t+1}\right)p^X_t &= -C^D_X\left(1+\zeta\Omega^D\right)\Big|_{t+1}+p^X_{t+1},
&&\text{hence}
&
p^X_t &= -\sum_{s=t+1}^\infty \Lambda_{t,s}\,C^D_X\left(1+\zeta\Omega^D\right)\Big|_s\;<\;0,
\label{eq:sp:costate:X}
\end{align}
with the compound discount factor $\Lambda_{t,s}$ of~\eqref{eq:sp:sum:discount}. $X$ is a pure liability: it accumulates the loss of the easy deposits, and its shadow price is the discounted stream of the extra exploration costs that loss imposes. It is what stops exploration from being a costless substitute for extraction in~\eqref{eq:sp:foc:D}.

Exhaustibility changes how this liability behaves in the long run, in a way worth recording. The divergence in~\eqref{eq:sp:setup:exhaustible-discovery} sits in the primal cost of discovering, not in the shadow price: $p^X$ never diverges. Indeed, once exploration has ceased for good it vanishes, because $C^D\left(0,X,t\right)=0$ for all $X$ and $t$ implies $C^D_X\left(0,X,t\right)=0$, so along any path with $D=0$ from period $t$ onward the dividend in~\eqref{eq:sp:costate:X} is identically zero and $p^X_t=0$. Past discoveries are a liability only insofar as exploration is still to come; when the ceiling has made further exploration worthless, the liability expires with it, and the corner condition~\eqref{eq:sp:sum:D-corner} thereafter compares $p^S$ alone against the diverging cost of the first tonne.

\smalltitle{Pollution.} A marginal tonne of the pollution stock lowers period-$t{+}1$ felicity by $v'\left(P\right)$ utils, lowers output by $\left|F_P\right|$, and decays. Its own transition passes on $\partial\left[P-\theta\left(P\right)P\right]/\partial P=1-\theta-\theta'P$ of the marginal tonne --- the \emph{marginal} retention factor. With $m^P=-\lambda p^P$,
\begin{align}
\left(1+r_{t+1}\right)p^P_t &= \underbrace{\mathcal V_{t+1}}_{\text{(b)}}+\underbrace{\left[1-\theta\left(P\right)-\theta'\left(P\right)P\right]_{t+1}}_{\text{(a)}}p^P_{t+1},
&
\mathcal V &\equiv \underbrace{\frac{v'\left(P\right)}{\lambda}}_{\text{utility damage}}-\underbrace{F_P\left(1-\zeta\Omega^Y\right)}_{\text{output damage}}\;>\;0,
\label{eq:sp:costate:P}
\end{align}
so that $p^P_t=\sum_{s>t}\Lambda_{t,s}\left[\prod_{j=t+1}^{s-1}\left(1-\theta-\theta'P\right)_j\right]\mathcal V_s>0$. The shadow cost of pollution is unambiguously positive: it is the discounted stream of marginal damages, both channels being costs. Term (b) is in final goods per tonne, the utility damage converted by $\lambda$ and the output damage scaled by $1-\zeta\Omega^Y$ for the same reason as everywhere else --- output lost to pollution also releases the embodied material it would have carried.

Term (a) is the effective survival factor, and it is worth marking that the decay entering it is the \emph{marginal} rate $\theta+\theta'P=\tfrac{d}{dP}\left[\theta\left(P\right)P\right]$ rather than the average rate $\theta$. Since $\theta'<0$ by~\eqref{eq:sp:setup:pollution-motion}, the marginal rate is the smaller of the two, so damages are discounted less heavily than the average decay rate would suggest. If decay saturates strongly enough that $\left(1-\theta-\theta'P\right)\ge1+r$ --- equivalently $r+\theta+\theta'P\le0$ --- the sum need not converge; that configuration is this model's version of a tipping point, and ruling it out is a restriction on $\theta\left(\cdot\right)$ rather than something the optimality conditions deliver.

\smalltitle{Embodied materials.} A marginal tonne of material embodied in the capital stock is released at rate $\delta$, at which point it becomes waste. It enters the period-$t{+}1$ terms of the Lagrangian~\eqref{eq:sp:lagrangian} in two places only, through $-\lambda p^W\delta M^K$ in the ledger term and through the surviving fraction $1-\delta$ of its own transition, so with $m^M=-\lambda p^M$ the costate condition reduces to
\begin{align}
\left(1+r_{t+1}\right)p^M_t &= \delta\,p^W_{t+1}+\left(1-\delta\right)p^M_{t+1},
&&\text{equivalently}
&
p^M_t &= \sum_{s=t+1}^\infty \Lambda_{t,s}\left(1-\delta\right)^{s-t-1}\delta\,p^W_s .
\label{eq:sp:costate:M}
\end{align}
The interpretation is exactly the postponement reading of term (b) of~\eqref{eq:sp:foc:Omega}: the tonne is released with probability-weight $\delta$ each period, each release costing $p^W$, and the stream is discounted both by $r$ and by the run-off of the liability itself.

In particular $p^M$ is a forward-weighted average of future values of $p^W$. It therefore shares the sign of $p^W$ whenever $p^W$ does not change sign from $t$ onward, but not otherwise. With $p^W$ and $r$ constant,
\begin{align}
p^M &= \frac{\delta}{r+\delta}\,p^W,
&
p^W-p^M &= \frac{r}{r+\delta}\,p^W .
\label{eq:sp:costate:M-steady}
\end{align}

\smalltitle{The waste stock.} A marginal tonne in the stockpile is drawn for handling with share $\mu$ in each period it survives. When handled, it absorbs $c^W$ final goods of collection and treatment, grossed up for their embodied material; a fraction $\alpha\varpi$ returns to production as secondary material, worth $\Psi-p^W$ per tonne net of the waste it will regenerate; and the remainder reaches the environment as the emission flow prices it. Since $\mathcal W$ acts on period $t+1$ only through $H=\mu\mathcal W$, every dividend term in $\partial\mathcal L^{t+1}/\partial\mathcal W_{t+1}$ is proportional to $\mu$, and collecting them together with the surviving fraction $1-\mu$ of the transition gives $\partial\mathcal L^{t+1}/\partial\mathcal W_{t+1}=\lambda_{t+1}\left[\mu h_{t+1}+\left(1-\mu\right)p^{\mathcal W}_{t+1}\right]$ with the \emph{handling dividend}
\begin{align}
h &\equiv
\underbrace{\alpha\varpi\left(\Psi-p^W\right)}_{\text{(a)}}
-\underbrace{c^W\left(1+\zeta\Omega^W\right)}_{\text{(b)}}
-\underbrace{p^P\left[1-\varpi\left(1-d^W\right)-d^W\alpha\varpi\right]}_{\text{(c)}} ,
\label{eq:sp:h}
\end{align}
in final goods per tonne, so that the costate condition for $m^{\mathcal W}=\lambda p^{\mathcal W}$ reads
\begin{align}
\left(1+r_{t+1}\right)p^{\mathcal W}_t &= \mu\,h_{t+1}+\left(1-\mu\right)p^{\mathcal W}_{t+1},
&&\text{hence}
&
p^{\mathcal W}_t &= \sum_{s=t+1}^\infty \Lambda_{t,s}\left(1-\mu\right)^{s-t-1}\mu\,h_s .
\label{eq:sp:costate:Wstock}
\end{align}
Term (a) of~\eqref{eq:sp:h} is the feedstock motive, carrying the same $\Psi-p^W$ as the treatment margin~\eqref{eq:sp:foc:varpi-gain} and for the same ledger reason: the recovered tonne must leave as waste all over again. Term (b) is the real resources absorbed in collection and treatment, $c^W\left(\varpi\right)=c^c\varpi+c^T\left(\varpi\right)$ being the average handling cost per handled tonne. Term (c) is the pollution damage, its bracket the fraction of the handled tonne that reaches the environment: $1$ at $\varpi=0$ and $d^W\left(1-\alpha\right)>0$ at $\varpi=1$, so a handled tonne always adds something to $\Xi$ however thoroughly it is processed. At the no-treatment corner the dividend reduces to $h=-p^P<0$: a stockpiled tonne in an economy that processes nothing is a pure pollution liability, and the stockpile is priced as one.

Three remarks. First, the forward representation makes the anthropogenic reserve the exact mirror of the virgin reserve: compare~\eqref{eq:sp:costate:S}. $S$ pays a dividend by making the extraction still to come cheaper; $\mathcal W$ pays one by delivering handled tonnes worth $h$ each, weighted by the geometric probability $\left(1-\mu\right)^{s-t-1}\mu$ that the tonne is handled in period $s$. On squeeze paths, where material values grow while handling costs stay bounded, $h>0$ and the stockpile is a genuine asset --- the same economics a model with same-period handling would express through a negative waste price, now attached to a stock that exists.

Second, the effective survival factor is larger than it looks. Substituting $p^W=-p^{\mathcal W}$ from~\eqref{eq:sp:foc:W} into term (a) and collecting the $p^{\mathcal W}_{t+1}$ terms,
\begin{align}
\left(1+r_{t+1}\right)p^{\mathcal W}_t &= \left[1-\mu\left(1-\alpha\varpi\right)\right]_{t+1}p^{\mathcal W}_{t+1}
+\mu\Big[\alpha\varpi\,\Psi-c^W\left(1+\zeta\Omega^W\right)-p^P\left[1-\varpi\left(1-d^W\right)-d^W\alpha\varpi\right]\Big]_{t+1},
\label{eq:sp:costate:Wstock-net}
\end{align}
so in marginal-value terms the stockpile runs off only at the \emph{leakage} share $\mu\left(1-\alpha\varpi\right)$ per period, not at the handling share $\mu$: the fraction $\alpha\varpi$ of each handled tonne re-enters production and eventually returns to the stockpile as waste, and the identification $p^W=-p^{\mathcal W}$ credits that return automatically.

Third, at a rest point of the price system the recursion~\eqref{eq:sp:costate:Wstock} solves to $p^{\mathcal W}=\mu h/\left(r+\mu\right)$ --- the handling dividend, discounted for the geometric wait until the tonne is handled --- and $p^W=-\mu h/\left(r+\mu\right)$ rearranges to
\begin{align}
p^W\left(1-\alpha\varpi+\frac{r}{\mu}\right) &= c^W\left(1+\zeta\Omega^W\right)+p^P\left[1-\varpi\left(1-d^W\right)-d^W\alpha\varpi\right]-\alpha\varpi\,\Psi ,
\label{eq:sp:costate:Wstock-limit}
\end{align}
which is, up to the correction term $r/\mu$, precisely the waste condition a model without a waste stock would impose in every period. The stock changes \emph{when} the waste price is paid --- one geometric draw at a time rather than at once --- and at a rest point that timing survives only as the $r/\mu$ term, which shrinks as handling gets faster relative to discounting. At the buffer corner $\mu=1$ the recursion instead collapses to $p^W_t=-h_{t+1}/\left(1+r_{t+1}\right)$: the full static margin, one period ahead and discounted.

\smalltitle{Consumption growth.} The rate $r$ is defined from the goods multiplier, but by~\eqref{eq:sp:foc:C} that multiplier is tied to marginal utility only up to the consumption wedge, $u'\left(C\right)=\lambda\left(1+\zeta\Omega^C\right)$. Taking the ratio of this condition across two adjacent periods,
\begin{align}
\frac{u'\left(C_t\right)}{\beta\,u'\left(C_{t+1}\right)} &= \left(1+r_{t+1}\right)\frac{1+\zeta_t\Omega^C_t}{1+\zeta_{t+1}\Omega^C_{t+1}},
&&\text{i.e.}
&
\eta\,\Delta\ln C_{t+1} &= \ln\frac{1+r_{t+1}}{1+\rho}-\Delta\ln\left(1+\zeta\Omega^C\right)_{t+1}
\label{eq:sp:keynes-ramsey}
\end{align}
where the log form holds under CRRA utility with $\eta\equiv-u''C/u'$ constant; for general $u$, the first equality is the operative condition. This is the Euler equation with one extra term. A consumption wedge that is constant over time drops out and affects only the level of consumption; it is a wedge that is \emph{changing} that tilts the path. A rising $\zeta\Omega^C$ --- material becoming scarcer relative to the storage margin --- depresses consumption growth below what the return on capital alone would justify, and a falling one raises it.
